% Pegar en el Overleaf del grupo, entre "Variables Creadas" y "Estadísticas Descriptivas".
% Requiere \usepackage{booktabs} en el preámbulo (si la tabla de descriptivas ya usa
% \toprule, ya está cargado).

% --- Lámina A: efecto esperado a priori de cada X ---------------------
\begin{frame}{Efecto esperado de cada variable sobre la $Y$ (a priori)}
  \centering
  \footnotesize
  \begin{tabular}{lcl}
    \toprule
    \textbf{Variable} & \textbf{Signo} & \textbf{Por qué} \\
    \midrule
    \texttt{PTJE\_NEM}, \texttt{PTJE\_RANKING} & $+$ & Pesan en el puntaje ponderado \\
    \texttt{CLEC}, \texttt{MATE1}               & $+$ & Más puntaje, más carreras alcanzables \\
    \texttt{BEA} (10\,\% mejor del curso)       & $+$ & Buen rendimiento y beca para pagar \\
    \texttt{PACE}                               & $+$ & Cupos garantizados en universidades PACE \\
    \midrule
    \texttt{INGRESO\_PERCAPITA}                 & $+$ & Puede pagar y preparar la PAES \\
    \texttt{INGRESO\_COLEGIO}                   & $+$ & Entorno con más recursos y redes \\
    \texttt{TP\_DUMMY}                          & $-$ & Menos preparación PAES; más salida a CFT \\
    \midrule
    \texttt{N\_MATRICULADOS}                    & $+$ & Cultura de seguir estudios en el colegio \\
    \texttt{N\_ALUMNOS}                         & ?   & Más recursos, pero menos atención por alumno \\
    \texttt{SEXO}                               & ?   & Mujeres: más NEM; hombres: más MATE1 \\
    \bottomrule
  \end{tabular}

  \medskip
  El ingreso podría pesar menos de lo esperado por la \textbf{gratuidad}:
  es parte de lo que el modelo tiene que medir.
\end{frame}

% --- Lámina B: variables que no pueden ser X ---------------------------
\begin{frame}{Las variables de la matrícula no pueden ser $X$}
  \centering
  \begin{tabular}{ll}
    \toprule
    \textbf{Variable} & \textbf{Problema} \\
    \midrule
    \texttt{PREFERENCIA}                  & Con ella se construye la $Y$ \\
    \texttt{PTJE\_POND}                   & Es el puntaje de la carrera en que ya quedó \\
    \texttt{VIA}, \texttt{TIPO\_MATRICULA} & Vacías en los 39.548 no matriculados \\
    \bottomrule
  \end{tabular}

  \medskip
  Usarlas sería explicar la $Y$ con la misma $Y$: quedan en la base solo
  como información.
\end{frame}

% --- Lámina C (respaldo, después de "¡Muchas Gracias!"): solo si preguntan por el N
\begin{frame}{Respaldo: ¿por qué no incluimos a quienes rindieron y no postularon?}
  \begin{itemize}
    \item La pregunta es \textbf{en qué preferencia quedó} el estudiante:
          quien no postuló no ordenó preferencias, así que no participa de
          ese evento.
    \item Sí incluimos a los \textbf{39.548 que postularon y no se
          matricularon}: tuvieron la opción y no quedaron arriba ($Y = 0$).
    \item Si se agregan los que no postularon, la $Y$ mezclaría dos
          decisiones distintas: \textbf{postular o no} y \textbf{dónde quedar}.
  \end{itemize}
\end{frame}

% --- En "Conclusiones y Próximos Pasos", cambiar el recuadro "Hacia la Etapa 3" por:
%   Hacia la Etapa 2: cruzar cada X con la Y y agregar variables del MINEDUC
%   por RBD (cursos, alumnos por profesor, ruralidad).
